\chapter{Structure of Config-File}
The basic structure of the config-file for CFS++ in our XML
syntax is as follows:

\begin{verbatim}
<?xml version="1.0"?>
<cfsSimulation xmlns="http://www.cfs++.org">

 ... Parameters ...

</cfsSimulation>
\end{verbatim}

The first line gives version information for the XML parser, while the
second line is the starting tag for the main enclosing element called
cfsSimulation.  The last line then is the closing tag for this main
element. All parameters must appear between the opening and the
closing tag.

Only a few tags are allowed in this main body (see table
\ref{tab:Layer1Elements}). They will be documented in the next
sections. These top level tags have to be used in the order shown in
table \ref{tab:Layer1Elements}.

\begin{table}[h!]
  \centering
  \begin{tabular}{lll}
    \hline\noalign{\smallskip}
    leave name & leave type & occurence constraints \\
    \noalign{\smallskip}\hline\noalign{\smallskip}
    analysis      & plane   & required\\
    geometry      & plane   & required\\
    input         & plane   & optional\\
    output        & plane   & optional\\
    materialData  & plane   & optional\\
    domain        & section & required \\
    pdeList       & section & required \\
    couplingList  & section & required for iterative or direct coupling\\
    multiSequence & section & required for multi sequence analysis \\
    transient     & section & required for transient analysis \\
    \noalign{\smallskip}\hline
  \end{tabular}
  \caption{\label{tab:Layer1Elements}Possible commands of the 1. level
    in the config-file}  
\end{table}



